\documentclass[10pt,a4paper]{amsart}
\usepackage{amsmath,amssymb,amsthm,graphicx}
\usepackage{amsmath}
\usepackage{amssymb}
\usepackage{hyperref}
\usepackage{amsmath,amssymb,amsthm,enumerate,mathabx}

 

\newtheorem{Theorem}{Theorem}[section]
\newtheorem{Proposition}[Theorem]{Proposition}
\newtheorem{Lemma}[Theorem]{Lemma}
\newtheorem{Corollary}[Theorem]{Corollary}
\newtheorem{Theorem1}{Theorem}
 
\theoremstyle{definition}
\newtheorem{Definition}[Theorem]{Definition}
\newtheorem{Fact}[Theorem]{Fact}
\newtheorem{Conjecture}[Theorem]{Conjecture}
\newtheorem{Example}[Theorem]{Example}
\newtheorem{Remark}[Theorem]{Remark}
\newtheorem{Question}{Question}
\newtheorem{Claim}{Claim}

\DeclareMathOperator{\cl}{cl_S}
\DeclareMathOperator{\dcl}{dcl}
\DeclareMathOperator{\tp}{tp}
\DeclareMathOperator{\Aut}{Aut}
\DeclareMathOperator{\conv}{conv} 
\DeclareMathOperator{\dom}{dom}
\DeclareMathOperator{\acc}{acc}

\usepackage[margin=23mm]{geometry}

\title{Vaught's conjecture for theories of discretely ordered structures}
\author{Predrag Tanovi\'c}
\thanks{This research was supported by the Science Fund of the Republic of 
Serbia, Grant No. 7750027: Set-theoretic, model-theoretic and 
Ramsey-theoretic phenomena in mathematical structures: similarity and 
diversity–SMART}
\address{Mathematical Institute SANU, Belgrade, Serbia} 
\email{tane [at] mi.sanu.ac.rs}
\title{Simple types cut out initial parts that are type-definable, discretely ordered and closed under minima}
\author{Predrag Tanović}
\date{arxiv:2212.13605; CC BY 4.0}
\begin{document}
\maketitle
\noindent\textit{Editorial note.} All delivered mathematical text is the author's own, copied from the pinned publisher source sources/vaught-discrete-orders-2212.13605v3/source0.tex; no statement or proof step was rewritten, reordered or removed.\par\medskip
\noindent\textit{Editorial note.} Preamble provenance: the retained preamble is the source's own layout adapter, if any, plus the author's macro block at lines 4-40 of sources/vaught-discrete-orders-2212.13605v3/source0.tex (everything between the publisher class metadata and \textbackslash{}begin\{document\}, comments excluded). Nothing was added, no macro was invented, and the publisher document class is deliberately not reproduced; the excerpt serves only as preamble, never as a span coordinate.\par\medskip
\noindent\textit{Editorial note.} External references are navigation pointers only; every deliberately kept step of the author's argument stays inside the counted spans. The displayed numbers of those pointers are the author's own numbering under the author's \textbackslash{}newtheorem declarations (the source has no subsections, so numbering is section.index), computed mechanically by the lane's scanner.\par\medskip
\noindent\textit{Editorial note.} The counted proof refers to its own clauses by name ('by part (c)') and to no labelled item of the paper, and it contains no bibliography citation, so the card needs neither a navigation stub nor a bibliography excerpt.\par\medskip
\noindent\textit{Editorial note.} The staged licence evidence for this source carries no status field, so the card points at a per-card reuse record generated inside its own task directory (license-reuse.json, material basis only).\par\medskip
\noindent\textit{Editorial note.} Source metadata: this source stages SOURCE\_RECORD.json (no SOURCE.json), which carries the title, the identifier and a list of source URLs, while the author list and the fixed-version URL live in the licence evidence; r001 was rejected with 'nonempty task field required: source\_url' before any compile. The generator now falls back across those files.\par\medskip
\noindent\textit{Allowed standard prerequisites (editorial).}\begin{itemize}
\item The definition of a simple type over A, witnessed by a subset C of the definable closure and a linear order (D,<). Set up in the definitions preceding the counted statement; used throughout the counted proof. Fixes the data the four clauses quantify over and the meaning of C-type used in each step.
\item The uniqueness of the simple type p extending the C-type condition. Used inside the counted proof at line 155. Lets a point whose type is finitely satisfied in C but not algebraic be identified with p, giving the initial-part claim.
\end{itemize}\medskip
\begin{Lemma}\label{Lemma_basic_simple}
If   $p\in S_1(A)$ is a simple type, as witnessed by $C\subset\dcl(A)$ and $(D,<)$, then:

(a) $P(M):=C\cup p(M)$ is an initial part  of $(D,<)$ that is type-definable over $A$.

(b) There exists an  $A$-definable, discretely ordered, initial part $D_0\subseteq D$  with $P(M)\subseteq D_0$. In particular, $(p(M),<)$ is a discrete order.  

(c) Every nonempty, relatively $M$-definable subset of  $P(M)$ has a  minimum.  

(d) If $C$ and $(D_1,<_1)$ also witness the simplicity of $p$ such that $<_1$ and $<$ agree on $C$, then they agree on $P(M)$, too. 
\end{Lemma}

\begin{proof}To simplify the notation, assume $A=\emptyset$ and  write $x<y$ instead of $x\in D\land y\in D\land x<y$. 

(a) Since $C$ is an initial part of $(D,<)$,    $C<p(M)$ holds. 
To prove that $P(M)$ is an initial part of $D$, it suffices to show that $a\in p(M)$ and $b< a$ together imply $b\in P(M)$. Let $\phi(x)\in \tp(b)$. 
Then  $\exists y(\phi(y)\land y< x)\in\tp(a)=p$. Since $p$ is   a $C$-type, there are $c\in C$ and $b'\in D$ with $\models\phi(b')\land b'< c$.
Then $b'< c$ implies $b'\in C$, so   $\phi(x)$ is realized by $b'\in C$. We have just shown that every formula $\phi(x)\in \tp(b)$ is realized in $C$, so the type $\tp(b)$ is finitely satisfied in $C$; 
If $\tp(b)$ is algebraic  then $b\in C$; otherwise, $\tp (b)$ is a $C$-type, so $\tp(b)=p$ holds by the uniqueness of $p$.
In both cases we have $b\in P(M)$, so $P(M)$ is an initial part of $D$. Then it is is easy to see that $P(M)$ is defined by  $P(x):=\{\exists y(\theta(y) \land   x<y)\mid \theta(y)\in p(y)\}$. 

\smallskip
(b) Let $c_0=\min C$ and let $\phi(x)$ be a formula expressing that the interval $[c_0,x]$ is discretely ordered by $<$.  For all $c\in C$ the interval $[c_0,c]$ is finite, so $\models  \phi(c)$ holds. Since $p$ is a $C$-type, $\phi(x)\in p$. Let $D_0$ be the union of all intervals $[c_0,a]\subset D$ where $a\in\phi(M)$.  

\smallskip
(c) Let $\emptyset\neq D_1\subset P(M)$  where $D_1$ is relatively defined by $\theta(x,\bar b)$ ($\bar b\in M$). Pick $a\in p(M)$ with $[c_0,a]\cap D_1\neq\emptyset$.  Let $\phi(x,\bar y)$ be a formula expressing that the set
$[c_0,x]\cap\theta(M,\bar y)$ is nonempty and has minimum. 
In particular, $\models\phi(a,\bar b)$ holds. For all $c\in C$ the interval $[c_0,c]$ is finite, so $\models \forall \bar y \left(\exists z(c_0\leq z\leq c\land \theta(z,\bar y))\rightarrow \phi(c,\bar y)\right)$. Since $p$ is a $C$-type and $a\in p(M)$, 
\begin{center} 
$\models \forall \bar y \left(\exists z(c_0\leq z\leq a\land \theta(z,\bar y))\rightarrow \phi(a,\bar y)\right)$.  
\end{center}
In particular,   $\models  \exists z(c_0\leq z\leq a\land \theta(z,\bar b))\rightarrow \phi(a,\bar b)$.  Now  $[c_0,a]\cap D_1\neq \emptyset$ implies $\models  \exists z(c_0\leq z\leq a\land \theta(z,\bar b))$, so $\models\phi(a,\bar b)$ and $\min ([c_0,a]\cap D_1)$ exists. Hence $\min D_1$ exists, too.

\smallskip
(d) Suppose not. Then for all $a\in p(M)$ the set $D(a)$ defined by $x<a\land x>_1 a$ is nonempty so, by part (c), $b=\min_{<} D(a)$ is well-defined. By part (a), $C<_1p(M)$ implies $b\notin C$, so  $b<a$ implies $b\in p(M)$. In particular, $D(b)\neq\emptyset$ and for any $b'\in D(b)$ we have $b'<b<a\land b'>_1b>_1 a$, contradicting the minimality of $b$ in $D(a)$. 
\end{proof}

\end{document}
